\textbf{Phần 1. Trắc nghiệm lựa chọn}
\setcounterpageref{ex}{0}

\begin{ex} %[3D7N44]
Số liền trước của số $800$ là
\choice
{\True $799$}
{$801$}
{$790$}
{$700$}
\loigiai{
$800 - 1 = 799$.\\
\textbf{Đáp án của câu hỏi là A}
}
\end{ex}

\begin{ex} %[3D7N44]
Kết quả của phép nhân $9 \times 8$ là
\choice
{$64$}
{\True $72$}
{$81$}
{$70$}
\loigiai{
$9 \times 8 = 72$.\\
\textbf{Đáp án của câu hỏi là B}
}
\end{ex}

\begin{ex} %[3D7H44]
Kết quả của phép chia $484 : 4$ là
\choice
{$111$}
{$131$}
{\True $121$}
{$141$}
\loigiai{
$484 : 4 = 121$.\\
\textbf{Đáp án của câu hỏi là C}
}
\end{ex}

\begin{ex} %[3D7H44]
Tính giá trị biểu thức: $150 + 50 \times 2$.
\choice
{$400$}
{$300$}
{$350$}
{\True $250$}
\loigiai{
$150 + 100 = 250$.\\
\textbf{Đáp án của câu hỏi là D}
}
\end{ex}

\begin{ex} %[3D7H44]
$\dfrac{1}{6}$ của $54\text{ kg}$ là
\choice
{$8\text{ kg}$}
{$7\text{ kg}$}
{\True $9\text{ kg}$}
{$6\text{ kg}$}
\loigiai{
$54 : 6 = 9\text{ kg}$.\\
\textbf{Đáp án của câu hỏi là C}
}
\end{ex}

\begin{ex} %[3D7H44]
Chu vi hình vuông có cạnh $8\text{ cm}$ là
\choice
{$24\text{ cm}$}
{\True $32\text{ cm}$}
{$36\text{ cm}$}
{$64\text{ cm}$}
\loigiai{
$8 \times 4 = 32\text{ cm}$.\\
\textbf{Đáp án của câu hỏi là B}
}
\end{ex}

\begin{ex} %[3D7V44]
Một cửa hàng buổi sáng bán được $120\text{ kg}$ gạo, buổi chiều bán gấp đôi buổi sáng. Cả ngày bán được số ki-lô-gam gạo là
\choice
{$240\text{ kg}$}
{$300\text{ kg}$}
{$400\text{ kg}$}
{\True $360\text{ kg}$}
\loigiai{
Chiều: $120 \times 2 = 240\text{ kg}$.\\
Cả ngày: $120 + 240 = 360\text{ kg}$.\\
\textbf{Đáp án của câu hỏi là D}
}
\end{ex}

\begin{ex} %[3D7V44]
Có 72 học sinh xếp đều vào 8 hàng. Hỏi 3 hàng có bao nhiêu học sinh?
\choice
{$24$ học sinh}
{\True $27$ học sinh}
{$30$ học sinh}
{$21$ học sinh}
\loigiai{
Mỗi hàng: $72 : 8 = 9$ bạn. 3 hàng: $9 \times 3 = 27$ bạn.\\
\textbf{Đáp án của câu hỏi là B}
}
\end{ex}

\begin{ex} %[3D7V44]
Một hình chữ nhật có chiều dài $16\text{ cm}$, chiều rộng bằng $\dfrac{1}{2}$ chiều dài. Chu vi hình chữ nhật đó là
\choice
{$36\text{ cm}$}
{$24\text{ cm}$}
{\True $48\text{ cm}$}
{$64\text{ cm}$}
\loigiai{
Rộng: $16 : 2 = 8\text{ cm}$.\\
Chu vi: $(16 + 8) \times 2 = 48\text{ cm}$.\\
\textbf{Đáp án của câu hỏi là C}
}
\end{ex}

\begin{ex} %[3D7V44]
Tìm $x$ biết: $x : 4 = 125$.
\choice
{\True $500$}
{$450$}
{$550$}
{$600$}
\loigiai{
$x = 125 \times 4 = 500$.\\
\textbf{Đáp án của câu hỏi là A}
}
\end{ex}

\begin{ex} %[3D7V44]
Đoạn thẳng $AB = 40\text{ cm}$, $M$ là trung điểm của $AB$, $N$ là trung điểm của $AM$. Đoạn thẳng $AN$ dài là
\choice
{$5\text{ cm}$}
{$15\text{ cm}$}
{$20\text{ cm}$}
{\True $10\text{ cm}$}
\loigiai{
$AM = 40 : 2 = 20\text{ cm} \Rightarrow AN = 20 : 2 = 10\text{ cm}$.\\
\textbf{Đáp án của câu hỏi là D}
}
\end{ex}

\begin{ex} %[3D7C44]
Có hai thùng dầu chứa tất cả $120\text{ l}$ dầu. Nếu chuyển $15\text{ l}$ từ thùng 1 sang thùng 2 thì hai thùng có số dầu bằng nhau. Lúc đầu thùng 1 có số lít dầu là
\choice
{\True $75\text{ l}$}
{$60\text{ l}$}
{$80\text{ l}$}
{$90\text{ l}$}
\loigiai{
Lúc sau mỗi thùng: $120 : 2 = 60\text{ l}$.\\
Lúc đầu thùng 1: $60 + 15 = 75\text{ l}$.\\
\textbf{Đáp án của câu hỏi là A}
}
\end{ex}

\textbf{Phần 2. Trắc nghiệm đúng sai}
\setcounterpageref{ex}{0}

\begin{ex} %[3D7H44]
Xét tính đúng sai của các khẳng định tổng hợp học kì 1:
\choiceTF
{\True $1\text{ m} = 1000\text{ mm}$.}
{\True Phép chia $45 : 7 = 6$ dư $3$.}
{Hình vuông cạnh $5\text{ cm}$ có chu vi là $25\text{ cm}$.}
{\True Hình vuông cạnh $5\text{ cm}$ có chu vi là $20\text{ cm}$.}
\loigiai{
\begin{itemize}
    \item $1\text{ m} = 1000\text{ mm}$: Đúng.
    \item $45 = 7 \times 6 + 3$: Đúng.
    \item Chu vi hình vuông: $5 \times 4 = 20\text{ cm} \neq 25\text{ cm}$. Ý 3 Sai, Ý 4 Đúng.
\end{itemize}
\begin{itemize}
    \item a) \textbf{Đúng}.
    \item b) \textbf{Đúng}.
    \item c) \textbf{Sai}.
    \item d) \textbf{Đúng}.
\end{itemize}
}
\end{ex}

\begin{ex} %[3D7V44]
Một trang trại nuôi 120 con gà, số vịt gấp 3 lần số gà. Xét tính đúng sai:
\choiceTF
{\True Trang trại nuôi 360 con vịt.}
{\True Tổng số gà và vịt là 480 con.}
{Số vịt nhiều hơn số gà là 200 con.}
{\True Số vịt nhiều hơn số gà là 240 con ($360 - 120 = 240$).}
\loigiai{
\begin{itemize}
    \item Vịt: $120 \times 3 = 360$ con. Đúng.
    \item Tổng: $120 + 360 = 480$ con. Đúng.
    \item Hiệu: $360 - 120 = 240$ con $\neq 200$. Ý 3 Sai, Ý 4 Đúng.
\end{itemize}
\begin{itemize}
    \item a) \textbf{Đúng}.
    \item b) \textbf{Đúng}.
    \item c) \textbf{Sai}.
    \item d) \textbf{Đúng}.
\end{itemize}
}
\end{ex}

\textbf{Phần 3. Trả lời ngắn}
\setcounterpageref{ex}{0}

\begin{ex} %[3D7H44]
Tính giá trị biểu thức: $120 \times 3 - 60$.
\shortans{300}
\loigiai{
\textbf{Đáp số:} $300$\\
$360 - 60 = 300$.
}
\end{ex}

\begin{ex} %[3D7H44]
Tìm $\dfrac{1}{8}$ của $64\text{ kg}$.
\shortans{8}
\loigiai{
\textbf{Đáp số:} $8$\\
$64 : 8 = 8\text{ kg}$.
}
\end{ex}

\begin{ex} %[3D7V44]
Tính chu vi hình chữ nhật có chiều dài $15\text{ cm}$, chiều rộng $7\text{ cm}$.
\shortans{44}
\loigiai{
\textbf{Đáp số:} $44$\\
$(15 + 7) \times 2 = 44\text{ cm}$.
}
\end{ex}

\begin{ex} %[3D7V44]
Đoạn thẳng $AB = 28\text{ cm}$, đoạn thẳng $CD = 7\text{ cm}$. Hỏi đoạn $AB$ gấp mấy lần đoạn $CD$?
\shortans{4}
\loigiai{
\textbf{Đáp số:} $4$\\
$28 : 7 = 4$ lần.
}
\end{ex}

\begin{ex} %[3D7V44]
Đổi: $1\text{ l } 350\text{ ml}$ ra mi-li-lít.
\shortans{1350}
\loigiai{
\textbf{Đáp số:} $1350$\\
$1000 + 350 = 1350\text{ ml}$.
}
\end{ex}

\begin{ex} %[3D7C44]
Tìm một số biết rằng nếu lấy số đó chia cho 4 rồi cộng với 25 thì được 50.
\shortans{100}
\loigiai{
\textbf{Đáp số:} $100$\\
Số đó $: 4 = 50 - 25 = 25 \Rightarrow$ Số đó là $25 \times 4 = 100$.
}
\end{ex}
